% ข้อมูลสำหรับแบบเสนอหัวข้องานวิจัย CS1
% แก้เฉพาะค่าด้านขวาของแต่ละคำสั่งก่อนส่งเอกสารจริง

\newcommand{\proposalcode}{CS1}
\newcommand{\proposalstudentid}{6611501658}
\newcommand{\proposalstudentname}{ทิวัตถ์ เลิศชัย}
\newcommand{\proposalstudyprogram}{ภาคในเวลา}
\newcommand{\proposalyearlevel}{4}
\newcommand{\proposalgroup}{วท.บ.661(4)/1}

\newcommand{\proposaltitleth}{การพัฒนาระบบจัดการและบันทึกการฝึกงานแบบมีส่วนร่วมหลายฝ่ายสำหรับคณะวิทยาศาสตร์ มหาวิทยาลัยราชภัฏจันทรเกษม}
\newcommand{\proposaltitleen}{Development of a Multi-stakeholder Internship Management and Logbook System for the Faculty of Science, Chandrakasem Rajabhat University}

\newcommand{\proposaladvisor}{ผู้ช่วยศาสตราจารย์วรรณา วิโรจน์แดนไทย}
\newcommand{\proposalchair}{รอตรวจสอบ}

% ช่วงเวลาด้านล่างเป็นแผนเบื้องต้น ให้แก้เป็นเดือนจริงก่อนส่ง
\newcommand{\proposalperiod}{8 เดือน นับจากวันที่หัวข้อได้รับอนุมัติ}

% เทคโนโลยีของระบบจริงยังไม่ปรากฏในเอกสารต้นทาง จึงตั้งค่าแบบไม่ผูกกับผลิตภัณฑ์
\newcommand{\proposalhardware}{คอมพิวเตอร์สำหรับพัฒนาและทดสอบระบบที่มีหน่วยประมวลผลแบบ 64 บิต หน่วยความจำอย่างน้อย 8 GB พื้นที่ว่างอย่างน้อย 20 GB อุปกรณ์พกพาที่มีเว็บเบราว์เซอร์และระบบระบุตำแหน่งสำหรับทดสอบการลงเวลา และเครื่องแม่ข่ายหรือบริการคลาวด์สำหรับติดตั้งระบบทดสอบ}
\newcommand{\proposalsoftware}{ระบบปฏิบัติการที่รองรับเครื่องมือพัฒนา เว็บเบราว์เซอร์รุ่นปัจจุบัน Git โปรแกรมแก้ไขซอร์สโค้ด Docker และ Docker Compose รวมถึงฐานข้อมูล เฟรมเวิร์กส่วนหน้า และเฟรมเวิร์กฝั่งเซิร์ฟเวอร์ตามรุ่นจริงของ Trainy ซึ่งต้องระบุชื่อและหมายเลขรุ่นก่อนส่งข้อเสนอ}
